\subsection{3.2 STEM 攀登}\label{stem-climb}

STEM 攀登是我们三项领域专属强化学习（RL）训练中时间最长的一次，旨在强化模型在单轮问题求解场景下的核心推理能力。它覆盖广泛的 STEM 领域，包括数学、物理、化学和竞赛编程。我们还观察到攀登期间 rollout 中反映出的模型行为演变，详见附录 D.1。由于数据质量是这次攀登成功的关键，我们重点介绍训练数据，以及我们为构建、验证和过滤这些数据而搭建的管线。

整个 STEM 攀登基于可验证数据对进行：为产生任务专属奖励 \(R _ { \mathrm { t a s k } } ( q , y _ { i } )\)，我们从 \(y _ { i }\) 中提取模型的最终答案，并将其与标准答案进行比较——既可使用 SymPy（Meurer 等，2017）之类的形式化验证器，也可使用 AI 评判器；对于竞赛编程，我们则将模型生成的代码片段在配套的题目专属测试用例集上运行。因此，每条 RL 数据实例要么是查询与标准答案组成的对 \(( q , a )\)，要么是查询与 n 个测试用例组成的对 \(( q , \{ t _ { 1 } , \ldots , t _ { n } \} )\)。在设计获取这类数据对的管线时，我们主要考虑三个标准：高质量、难度适当和主题多样。

借助我们的 STEM 数据管线，我们处理了数百万份文档，产出用于 STEM 攀登的 STEM Mix 数据集，包含超过 5M 个样本；其中最具挑战性的部分包含超过 550k 对 (q, a) 数据。该数据集的一些特性见图 16。



\pandocbounded{\includegraphics[keepaspectratio,alt={image}]{images/7a68d3c760cdbe6cf1d3ea668ea712afa02216898685011fdd6c578d8820a4af.jpg}}

图 17. 从教科书和学术 PDF 中提取 (q, a) 数据对以构建 STEM Mix 的管线。

\subsection{3.2.1 数据管线}\label{data-pipeline}

我们的 STEM 数据管线将异构的原始来源——教科书、学术 PDF、论坛讨论、竞赛档案，以及从不同供应商获取的题目——转化为 (q, a) 数据对数据集。该管线基于可组合的、分阶段的架构构建：每个处理步骤都实现为独立的异步阶段。管线的来源专属实例按需选择和组合这些阶段中的子集。

下面我们介绍若干管线阶段，按四个阶段分组，以图 17 所示从长篇教科书和学术 PDF 中提取 (q, a) 数据对的示例为参照。对噪声或幻觉敏感的阶段会多次运行并采用共识投票。

分层解析。该阶段包含将原始文档转换为初始 (q, a) 数据对的各个步骤。根据文档类型，OCR 通过视觉-语言模型或 OCR 服务完成。文档被分块，无 STEM 内容的页面被丢弃，并移除样板文本。单独的阶段利用 OCR 输出中的结构线索构建分层表示，并修复断裂的交叉引用、错位的编号以及跨页拆分造成的伪影。最后，在问答（QA）抽取阶段，LLM 在清理后的文档结构中标记问题和答案片段，生成候选 (q, a) 数据对。

QA 配对。对于问题和答案出现在不同位置的数据源（例如章节末尾的习题、附录中的答案），多轮配对阶段将每个问题与其对应的答案进行匹配。问题首先利用结构信号和语义相似度与候选答案匹配；随后 LLM 从检索到的候选中选择最佳匹配的答案并验证配对。

整理。整理阶段的各步骤对 (q, a) 数据对进行标注、过滤和改写，以确保质量和格式一致性。其中若干步骤是基于 LLM 的分类：

• 可验证性：条目被分为可验证与不可验证两类；不可验证的条目（例如没有可核查答案的开放式论坛讨论）会被丢弃。

• 题目类型：我们区分不同的题目类型，包括开放式问题、选择题和证明题。

• 分类体系：分层主题分类器为题目分配细粒度的 STEM 主题，涵盖数学、物理、化学、生物学、计算机科学和工程学。

• PII 检测：专门的分类器标记个人身份信息（PII）；被标记的条目会被丢弃。



• 答案泄漏：检测并丢弃答案 a 明显包含在问题 q 中的 (q, a) 数据对。

另外两个整理阶段会对题目进行转换，使其对 RL 攀登更有用。转换（Conversion）将选择题和证明题改写为开放式形式；选择题往往可以通过猜测解决，会提供不可靠的奖励信号，而证明题不依赖更强的 AI 评判器则难以直接验证。该步骤运行三次，随后进行共识判定；未达成共识的条目被视为无法可靠转换而被丢弃。清理（Cleanup）移除多余的、非数学的文本以及与题目无关的外部引用。

评分。这一最终阶段量化题目难度，并过滤掉标准答案可能有误的条目。在求解阶段，四个模型档位各对每个问题求解 k 次，各档位的能力以其 AIME 2025 成绩作为代理指标。通过率被用于将条目归入不同的难度区间。

盲评阶段用于防范错误的标准答案。对于最强档位求解通过率较低的条目，我们将该模型的共识答案与标准答案以随机顺序呈现给评判器。如果评判器更倾向于共识答案，我们就判定该条目标准答案可疑并将其丢弃；如果它更倾向于标准答案，则将其保留为真正困难的问题。

\subsection{3.2.2 竞赛编程数据}\label{competitive-coding-data}

竞赛编程数据使用一条独立的专用管线，与我们的主 STEM 管线有所不同；这是因为竞赛编程题目需要为每个提示提供一组测试用例 \(t _ { 1 } , \ldots , t _ { n }\)。在 PDF 文档等非结构化来源中，很难找到带有全面测试用例的提示。因此我们依赖定向来源和供应商获取的数据，这省去了 STEM 数据管线中的许多提取和过滤步骤。

对于每个题目，我们获取参考解法并验证其能通过所有相关测试用例。我们共从多个来源收集了 160k 个题目，覆盖多种主题，包括分治、动态规划、图与树算法以及搜索算法。除测试用例之外，每个题目还包含运行时间和内存限制。我们的最终数据集支持 17 种编程语言，包括 Python、C++、C\#、Java、JavaScript、Rust 和 TypeScript。

\subsection{3.2.3 去重与去污染}\label{deduplication-and-decontamination}

我们的 STEM Mix 数据集和竞赛编程数据集都做了自身去重，并与我们在第 H 节报告的基准测试以及我们内部跟踪的奥林匹克竞赛和研究生级别 STEM 评测进行去重。我们使用三阶段管线进行去重：

\begin{enumerate}
\def\labelenumi{\arabic{enumi}.}
\item
  精确去重。作为第一步，我们使用 SHA-256 题目哈希识别完全重复的题目。
\item
  词法模糊去重。我们使用字符级 n-gram shingling 结合 MinHash 局部敏感哈希来识别近似重复的题目。Jaccard 相似度高于给定阈值的题目对会被标记为重复。
\item
  向量去重。最后，使用轻量级嵌入模型对题目进行嵌入。余弦相似度高于给定阈值的题目对会被标记为重复。
\end{enumerate}

我们在每个阶段仔细调优超参数，以便在确保严格排除任何基准泄漏的同时，尽可能多地保留数据。
